\chapter{Индексный дефект и метрика редуцированного связывающего quotient}
\label{ch:v7-index-defect-reduced-linking-quotient-gate}

\section{Почему локальный проход ещё нужно нормировать}

Предыдущий гейт получил правильный гессиан из однократного полярного чтения
\begin{equation}
 Z=T_U(C_s,C_t)=\frac12(UC_sU^*+C_t),
 \qquad
 \mathcal S_{\rm raw}=\frac12\|Z\|_{\rm HS}^2.
 \label{eq:v7-reduced-raw-quotient-action}
\end{equation}
Но коэффициент перед нормой нельзя переносить через quotient без проверки.
Исходное физическое действие уже имело вид половины суммы двух углов:
\begin{equation}
 \mathcal S_V
 =\frac12\left(\|C_s\|_{\rm HS}^2+\|C_t\|_{\rm HS}^2\right).
 \label{eq:v7-reduced-original-two-corner-action}
\end{equation}

На согласованных десятимерных опорах положим
\begin{equation}
 X=UP C_sPU^*,
 \qquad Y=C_t,
 \qquad
 Z=\frac{X+Y}{2},
 \qquad D=\frac{X-Y}{2}.
 \label{eq:v7-reduced-symmetric-antisymmetric-curvature}
\end{equation}
Тогда параллелограммное тождество даёт
\begin{equation}
 \frac12\left(\|X\|_{\rm HS}^2+\|Y\|_{\rm HS}^2\right)
 =\|Z\|_{\rm HS}^2+\|D\|_{\rm HS}^2.
 \label{eq:v7-reduced-parallelogram-identity}
\end{equation}
Следовательно, метрика, индуцированная исходным двухугловым носителем на
диагональном quotient, равна $\|Z\|^2$, а не $\frac12\|Z\|^2$.

\section{Нулевой гессиан индуцированной метрики}

В нуле согласованные фоновые кривизны совпадают после полярного переноса:
\begin{equation}
 U(A_0^*A_0)U^*=A_0A_0^*.
 \label{eq:v7-reduced-background-intertwining}
\end{equation}
Антисимметричная часть $D$ и все блоки ранга-один дефекта начинаются с
квадратичного порядка по полю. Поэтому они не дают вклада в гессиан нуля,
и машинный аудит получает точное тождество
\begin{equation}
 \Hess_0\|Z\|_{\rm HS}^2
 =\Hess_0\mathcal S_V.
 \label{eq:v7-reduced-inherited-hessian-identity}
\end{equation}

Удобно записать семейство метрик
\begin{equation}
 \mathcal S_c=\mathcal S_E+\frac c2\|Z\|_{\rm HS}^2.
 \label{eq:v7-reduced-metric-scale-family}
\end{equation}
Все двадцать тяжёлых направлений положительны точно при
\begin{equation}
 0\le c<\frac{16}{15}.
 \label{eq:v7-reduced-metric-scale-window}
\end{equation}
Сырая однокопийная конвенция $c=1$ проходит, но индуцированная метрика
$c=2$ возвращает прежний равный вес:
\begin{align}
 c=1:&\quad (n_-,n_0,n_+)=(7,0,20),
 &\lambda_{\rm heavy}^{\min}&=\frac25,
 \label{eq:v7-reduced-raw-pass}\\
 c=2:&\quad (n_-,n_0,n_+)=(21,0,6),
 &\lambda_{\rm heavy}^{\min}&=-\frac{28}{5}.
 \label{eq:v7-reduced-inherited-fail}
\end{align}
Обе метрики сохраняют положительный гессиан уже выбранного вакуума. Различие
касается именно эндогенного запуска из нуля.

\section{Проверка накопленных механизмов редукции}

\subsection{Real-полуслед}

Полное Real-удвоение действует одновременно на рёберный и физический
Hodge-блоки. Предыдущий строгий гейт уже доказал, что один физический
полуслед сокращает оба удвоения и сохраняет относительный вес один. Он не
превращает индуцированную метрику $c=2$ в $c=1$ только для физического
quotient.

\subsection{Лесной семейный quotient}

Лесной quotient главы~\ref{ch:v7-real-arrow-bimodule-forest-quotient-gate}
действует на вершинные семейные кадры после выбора ненулевого вакуума. Он
уменьшает $54$ ориентационных нуля до девяти циклических мод, но не
отождествляет source/target Gram-углы и не меняет их Hodge-метрику.

\subsection{Представленный junk}

На вспомогательном Real-модуле стрелок обычный модуль внутренних одноформ
равен нулю. Поэтому в текущем представленном исчислении нет ненулевого
идеала форм, факторизация по которому могла бы породить именно карту
$C_s\oplus C_t\mapsto Z$ вместе с новой нормой.

\subsection{BRST--BV и gauge-quotient}

Для вершинного калибровочного действия
\begin{equation}
 \delta A=\xi_tA-A\xi_s
 \label{eq:v7-reduced-gauge-variation}
\end{equation}
в точке $A=0$ получаем $\delta A=0$. Поэтому ненулевая отрицательная мода
гессиана нуля не является вертикальной gauge-модой и не может исчезнуть
только после BRST--BV-факторирования.

\begin{theorem}[No-go наследуемой метрики текущих редукций]
Полярный quotient канонически определяет общий десятимерный образ
кривизн, но метрика, наследуемая от исходного двухуглового Hodge-носителя,
считает его с коэффициентом $c=2$. Она даёт сигнатуру $(21,0,6)$.
Ни общий Real-полуслед, ни прежний лесной quotient, ни текущий junk-фактор,
ни gauge/BRST--BV-редукция не выводят сырую метрику $c=1$, обеспечившую
локальный проход предыдущего гейта.
\end{theorem}

\begin{nogocriterion}[Запрет потери половины при смене носителя]
Нельзя после отождествления двух согласованных углов сохранять коэффициент
$1/2$ перед однокопийным следом только потому, что он даёт правильную
сигнатуру. Требуется либо функториально вывести quotient-меру с этой
нормировкой, либо получить эквивалентное ослабление из смешанной
полярно-передаточной кривизны одного большего оператора.
\end{nogocriterion}

\begin{protocol}
\begin{enumerate}
 \item полярный carrier $10+10+1$: \textbf{сохранён};
 \item сырой quotient $c=1$: \textbf{локально проходит};
 \item индуцированный quotient $c=2$: \textbf{проваливает запуск};
 \item критическая граница: \textbf{$c=16/15$};
 \item Real-полуслед как источник $c=1$: \textbf{нет};
 \item лесной семейный quotient как источник $c=1$: \textbf{нет};
 \item текущий junk-фактор как источник $c=1$: \textbf{нет};
 \item BRST--BV как удаление лишней моды нуля: \textbf{нет};
 \item физическая quotient-метрика: \textbf{не выведена};
 \item следующий гейт: \textbf{смешанная полярно-передаточная кривизна}.
\end{enumerate}
\end{protocol}

Машинный сертификат сохранён в
\path{s2t_v7_index_defect_reduced_linking_quotient_gate_results.json}.

\begin{thebibliography}{9}
\bibitem{v7-reduced-weiner}
M.~Weiner,
\emph{On Quasi-Orthogonal Systems of Matrix Algebras},
arXiv:1002.0017.
\bibitem{v7-reduced-kodaka}
K.~Kodaka, T.~Teruya,
\emph{The Strong Morita Equivalence for Inclusions of $C^*$-Algebras and
Conditional Expectations for Equivalence Bimodules},
arXiv:1609.08263.
\end{thebibliography}